\subsection{Dataset Description}
\label{sec:dataset}

Training and evaluation rely on a paired laser ultrasonic dataset acquired under controlled conditions. The dataset follows a paired acquisition protocol: input signals are recorded with limited averaging (N shots per location), while reference targets use high averaging (M shots per location, with M \textgreater{}\textgreater{} N) to approximate ground-truth clean signals. Representative values of N = 16 and M = 256 are used throughout this study, yielding a nominal SNR difference of approximately 12 dB between input and target; however, the exact N and M values are determined by the experimental hardware and acquisition time constraints.

\textbf{Training and validation specimens.} All paired training and validation data are collected from an intact aluminum plate (300 mm \times{} 300 mm \times{} 3 mm). Aluminum is selected as the specimen material because it is a common aerospace structural alloy with well-characterized Lamb wave propagation behavior at ultrasonic frequencies, enabling reliable interpretation of acoustic features such as arrival times and waveform morphology. The intact plate provides a controlled environment for establishing paired data with consistent noise characteristics across multiple acquisition sessions.

\textbf{Defect test specimens.} Generalization performance is evaluated on aluminum specimens containing machined defects including through-thickness slots (width: 1.0 mm; length: 20 mm) and edge notches (depth: 1.5 mm; width: 0.5 mm). These machined defects produce measurable reflections and mode conversions in the Lamb wave response, making them suitable for assessing whether denoising preserves diagnostically relevant features. Separate defect-containing specimens are used exclusively for testing --- no defect data are used during training or validation.

\textbf{Data splitting.} The intact plate dataset is split 80/10/10 by acquisition session, yielding approximately 80\% of paired signals for training, 10\% for validation (used for early stopping and hyperparameter tuning), and 10\% for held-out testing of intact-plate performance. The defect test set consists of entirely separate specimens and is not part of the splitting ratio. This ensures that the test set provides an unbiased evaluation of cross-domain generalization from intact training data to defect-containing scenarios.

\textbf{Signal preprocessing.} All acquired signals are normalized to unit variance and windowed to isolate the primary Lamb wave arrival within a 50 \textmu{}s gate. Signals with observable acquisition artifacts (e.g., laser power fluctuations or probe misalignment) are discarded prior to dataset assembly, yielding a clean training corpus of approximately 10,000 paired signal instances across all sessions.
